\documentclass[11pt,a4paper]{article}
\usepackage[italian]{babel}
\usepackage[margin=2.5cm]{geometry}
\usepackage{parskip}
\usepackage{amsmath, amssymb, amsthm}
\usepackage{booktabs}
\usepackage{array}
\usepackage{xcolor}
\usepackage{needspace}
\usepackage{enumitem}
\usepackage{listings}
\usepackage{tikz}
\usepackage[most]{tcolorbox}
\usepackage[hidelinks]{hyperref}
\usetikzlibrary{decorations.pathreplacing, shapes.geometric, arrows.meta, positioning}

% Niente vedove/orfane: i paragrafi non lasciano righe isolate a fine/inizio pagina
\widowpenalty=10000
\clubpenalty=10000

% Elenchi più compatti
\setlist{itemsep=2pt, parsep=0pt, topsep=4pt, partopsep=0pt}

% --- Riquadri con bordo nero, senza numero (si spezzano tra due pagine) ---
% Uso: \begin{definizione} ... \end{definizione}
%      \begin{teorema}[Nome] ... \end{teorema}  (il nome è facoltativo)
\tcbset{nota/.style={enhanced jigsaw, breakable, colback=white,
  colframe=black, boxrule=0.5pt, sharp corners}}
\NewTColorBox{definizione}{}{nota,
  before upper={\textbf{Definizione.}\ }}
\NewTColorBox{teorema}{o}{nota,
  before upper={\textbf{Teorema\IfValueT{#1}{ (#1)}.}\ }}
\NewTColorBox{proposizione}{o}{nota,
  before upper={\textbf{Proposizione\IfValueT{#1}{ (#1)}.}\ }}
\NewTColorBox{proprieta}{o}{nota,
  before upper={\textbf{Proprietà\IfValueT{#1}{ (#1)}.}\ }}
\NewTColorBox{assioma}{o}{nota, boxrule=1pt,
  before upper={\textbf{Assioma\IfValueT{#1}{ (#1)}.}\ }}

% --- Ambienti semplici (senza riquadro) ---
% Il titolo sta in cima al blocco, anche se il contenuto inizia con un riquadro o
% con codice; e non resta mai da solo in fondo a una pagina.
% Uso: \begin{esempio}[Titolo facoltativo] ... \end{esempio}
\newcommand{\newrunin}[2]{%
  \NewDocumentEnvironment{#1}{o}{%
    \par\Needspace{4\baselineskip}\addvspace{\medskipamount}\noindent
    \textbf{#2}\IfValueT{##1}{ (##1)}\textbf{.}\ \ignorespaces}%
  {\par\addvspace{\medskipamount}}}
\newrunin{esempio}{Esempio}
\newrunin{esercizio}{Esercizio}
\newrunin{osservazione}{Osservazione}

% V e F nelle tavole di verità
\newcommand{\V}{\textsf{V}}
\newcommand{\F}{\textsf{F}}

% --- Codice C ---
% Uso: \begin{codice} ... \end{codice}      codice C numerato
%      \begin{comando} ... \end{comando}    comandi da terminale
%      \begin{risultato} ... \end{risultato} output di un programma
%      \lstinline|printf("ciao");|           codice dentro al testo
\lstdefinestyle{c}{
  language=C,
  basicstyle=\ttfamily\small,
  keywordstyle=\color{blue}\bfseries,
  commentstyle=\color{gray}\itshape,
  stringstyle=\color{red!70!black},
  numbers=left,
  numberstyle=\tiny\color{gray},
  numbersep=8pt,
  columns=flexible,
  keepspaces=true,
  breaklines=true,
  showstringspaces=false,
  tabsize=4,
  morekeywords={include, define},
  % lettere accentate nei commenti
  literate={à}{{\`a}}1 {è}{{\`e}}1 {é}{{\'e}}1 {ì}{{\`i}}1
           {ò}{{\`o}}1 {ù}{{\`u}}1 {È}{{\`E}}1 {À}{{\`A}}1
}
\lstset{style=c}
\newtcblisting{codice}{listing only, nota, left=6mm,
  listing options={style=c}}
\newtcblisting{comando}{listing only, nota,
  listing options={basicstyle=\ttfamily\small, numbers=none, language={},
    columns=flexible, keepspaces=true, breaklines=true}}
\newtcblisting{risultato}{listing only, enhanced jigsaw, breakable,
  colback=black!5, colframe=black, boxrule=0.5pt, sharp corners,
  listing options={basicstyle=\ttfamily\small, numbers=none, language={},
    columns=flexible, keepspaces=true, breaklines=true}}

% --- Diagrammi di flusso (simboli standard) ---
\tikzset{
  terminale/.style = {ellipse, draw, minimum width=2.5cm, minimum height=0.9cm},
  io/.style        = {trapezium, trapezium left angle=70, trapezium right angle=110,
                      draw, minimum width=2.5cm, minimum height=0.9cm},
  processo/.style  = {rectangle, draw, minimum width=2.5cm, minimum height=0.9cm},
  decisione/.style = {diamond, draw, aspect=2, minimum width=2.5cm},
  freccia/.style   = {thick, -{Stealth}}
}

\title{Immagine e controimmagine}
\author{Marco Cerbella}
\date{Lezione 4 -- 5 ottobre 2026}

\begin{document}
\maketitle

\section{Immagine e controimmagine}

Sia $f : D \to \mathbb{R}$ e sia $A$ un sottoinsieme di $D$.

\begin{definizione}
L'\emph{immagine di $A$ attraverso $f$} è l'insieme di tutte le immagini degli elementi di $A$:
\[
f(A) = \{f(x) \mid x \in A\} \subseteq \mathrm{Im}(f).
\]
Per $A = D$ si ha $f(D) = \{f(x) \mid x \in D\} \subseteq \mathbb{R}$, l'\emph{immagine di $f$}, $\mathrm{Im}(f)$.
\end{definizione}

\begin{definizione}
Sia $y \in \mathbb{R}$. La \emph{controimmagine} di $y$ attraverso $f$ è l'insieme
\[
f^{-1}(y) = \{x \in D : f(x) = y\}.
\]
Tale insieme è vuoto se e solo se $y \notin \mathrm{Im}(f)$. Sia $B \subseteq \mathbb{R}$: la \emph{controimmagine} di $B$ è
\[
f^{-1}(B) = \{x \in D : f(x) \in B\},
\]
cioè l'unione di tutte le controimmagini degli elementi di $B$.
\end{definizione}

\section{Esempi}

\begin{esempio}[$f(x) = x^2$]
Sia $f : \mathbb{R} \to \mathbb{R}$, $f(x) = x^2$. Allora
\begin{itemize}
    \item $f([1, 2]) = [1, 4]$;
    \item $f^{-1}([1, 4]) = [-2, -1] \cup [1, 2]$;
    \item $f^{-1}([-1, 0[) = \emptyset$, perché $[-1, 0[ \,\cap\, \mathrm{Im}(f) = \emptyset$;
    \item $f(\mathbb{R}) = [0, +\infty)$.
\end{itemize}
\end{esempio}

\begin{esempio}[Funzione segno]
Sia $g : \mathbb{R} \to \mathbb{R}$,
\[
g(x) = \operatorname{sign}(x) = \begin{cases} 1 & \text{se } x > 0, \\ 0 & \text{se } x = 0, \\ -1 & \text{se } x < 0. \end{cases}
\]
Allora $g([2, 3]) = \{1\}$; $g^{-1}(1) = \,]0, +\infty[$; $g^{-1}(3) = \emptyset$ (perché $3 \notin \mathrm{Im}(g)$); $g(\mathbb{R}) = \{-1, 0, 1\}$.
\end{esempio}

\begin{esempio}[Parte intera]
Sia $h : \mathbb{R} \to \mathbb{R}$, $h(x) = [x]$, la \emph{parte intera} di $x$, cioè il più grande intero $\leq x$. Allora
\[
h([0, 1[) = \{0\}, \quad h([1, 2[) = \{1\}, \quad h([-2, -1[) = \{-2\},
\]
$h(\mathbb{R}) = \mathbb{Z}$, e $h^{-1}(\tfrac12) = \emptyset$ perché $\frac12 \notin \mathbb{Z}$.
\end{esempio}

\begin{esempio}[Mantissa]
Sia $\ell : \mathbb{R} \to \mathbb{R}$, $\ell(x) = x - [x]$, la \emph{mantissa}. Per $x \in [n, n+1[$, con $n \in \mathbb{Z}$, si ha $\ell(x) = x - n$: il grafico è formato da infiniti segmenti paralleli di pendenza $1$, e $\ell(\mathbb{R}) = [0, 1[$. Ad esempio
\[
\ell(\tfrac12) = \tfrac12 - 0 = \tfrac12, \qquad
\ell(-\tfrac14) = -\tfrac14 - [-\tfrac14] = -\tfrac14 + 1 = \tfrac34.
\]
Inoltre
\[
\ell\big([-\tfrac14, \tfrac14]\big) = [\tfrac34, 1[ \,\cup\, [0, \tfrac14],
\qquad
\ell^{-1}\big([0, \tfrac12[\big) = \bigcup_{n \in \mathbb{Z}} \,[n, n + \tfrac12[.
\]
\end{esempio}

\end{document}
